% ============================================================================
% Section 7 -- related work. Claim C4 and the novelty argument.
%
% Source: seven papers read end-to-end 2026-08-10.
% CUT TO ~0.75pp 2026-09-27. Only the two papers a reviewer would reject on get
% a paragraph here: DirectKV (the competing mechanism) and Pie (the closest
% prior placement). NanoFlow's attack on the premise, the migration and
% KV-shrinking alternatives, NEO's compute axis and the endpoint argument are all
% in appendix C, which exists for exactly this.
%
% PROVENANCE: every number is attributed to the paper that reported it, which is
% stronger than a % src: into our own data.
% ============================================================================
\section{Related work}
\label{sec:related}

\paragraph{The mechanism is solved; DirectKV is a competitor, not a substrate.}
A trace-driven study of KV offloading finds device memory capacity, not token
budget, limits batch size~\cite{bottlenecks} and names scheduler improvements as
future work --- the premise this paper builds on.
DirectKV~\cite{directkv} then contributes zero-copy GPU attention over
CPU-resident KV on coherent interconnects, eliminating the staging buffer. Its
kernel reads host memory \emph{during} the attention that needs it, where we
prefetch layer $\ell{+}1$ while layer $\ell$ computes: they pay no device memory
and get no overlap, we pay one $k$-window and hide about $90\%$ of the transfer.
Neither dominates on paper and settling it needs both on one machine, which needs
their SM90a-only kernel. Two things their paper establishes matter here. Its
\S8.1 states the policy dimension is untouched --- the design ``does not change
scheduling, batching, or KV-management policies'' --- and its own analysis reaches
our PCIe conclusion from the kernel side, ``on PCIe, DirectKV should be viewed
primarily as a capacity-extension mechanism,'' while reporting NVLink-C2C cutting
attention latency up to $4.2\times$. Neither their evaluation nor ours covers a
70B-class configuration, which our model identifies as the most favourable point
of its curve, so that point is measured by nobody yet.

\paragraph{The closest prior mechanism partitions by layer.} Pie~\cite{pie}
targets the same hardware, stages KV through GPU buffers as we do, prefetches per
layer so transfer overlaps compute, and reports up to $1.9\times$ over vLLM.
\textbf{The difference is the axis of partition.} Pie relocates $m$ of $n$
\emph{layers} of KV, ``using per-layer KV cache as the basic management unit'',
and that split applies to \emph{every request identically} under one global
expansion factor. So \textbf{no request's KV is ever fully absent from the
device}, and \textbf{interconnect bandwidth is spent uniformly} --- a fixed
fraction of every request's KV rather than the whole KV of requests we select,
which is what makes \emph{which requests} a decision at all. The axes are not
orthogonal: they spend the same link budget and in the fluid limit reach the same
interior point, so our case for requests rests on what that limit omits
(\S\ref{sec:intro}). One thing favours layers and we concede it --- a layer split
has no integer quantization, so at long context it reaches operating points a
request split cannot express.

\paragraph{What else could be done instead.} Three alternatives bound the
contribution and Appendix~\ref{app:related} treats each in full.
NanoFlow~\cite{nanoflow} argues the opposite premise, that serving is
\emph{compute}-bound once operations of different resource classes are overlapped
within a device; the claims collide only above the batch a capacity limit permits,
and at batch 142 for an 8B model at 4K the step is still memory-bound by more than
$6\times$ (Appendix~\ref{app:params}). Llumnix~\cite{llumnix} live-migrates requests
\emph{across} instances instead of moving KV down a tier --- an economic choice
rather than a technical one, since migration needs a second instance with room while
tiering needs host DRAM already installed and idle. And architectures that shrink
the cache, MLA~\cite{deepseekv2} and sliding-window attention, do not by themselves
close the gate: by Eq.~\ref{eq:cancel} a uniform reduction in $\btok$ scales both
cohorts and leaves their ratio unchanged. We regard a stack built entirely on
MLA-class models as outside this paper's regime.

\paragraph{Summary of the gap.} Table~\ref{tab:axes} places the systems sharing our
hardware assumption on the two questions that matter. What is unsolved is which
requests live in which tier, how many a step can absorb without breaking the resident
cohort's latency, and when a request should move --- DirectKV states as much
explicitly, Fang et al.~\cite{calkv} close by calling for runtime policies, and
OmniServe~\cite{omniserve} takes its service classes as externally supplied.

% Single-column (2026-09-27): spanning both columns forced this to the top of a
% page and left a gap when it could not, which four rows do not earn.
\begin{table}[t]
\centering
\footnotesize
\setlength{\tabcolsep}{4pt}
\begin{tabular}{@{}llll@{}}
\toprule
 & placement unit & decided by & device KV \\
\midrule
vLLM, SGLang & --- & --- & all of it \\
Pie~\cite{pie} & layer, uniform & global factor & $(n{-}m)/n$ each \\
DirectKV~\cite{directkv} & none (all host) & nothing & none \\
\textbf{this paper} & \textbf{request} & \textbf{3 policies} & cohort only \\
\bottomrule
\end{tabular}
\caption{Systems either place nothing (both endpoints) or a fixed fraction of
everything (Pie); none decides per request. Our three policies are placement,
admission and migration.}
\label{tab:axes}
\end{table}
